%! Linear Regression — Unit 2, Lesson 2.5 — Slides (Beamer)
\documentclass[aspectratio=169, 11pt]{beamer}
\usepackage{algebra2-beamer}   % provides \forestheader{} and \sectionlabel[color]{}
% Coordinate grid: {xmin}{xmax}{ymin}{ymax}
\newcommand{\lgrid}[4]{%
  \draw[step=1, linegray!40, very thin] (#1,#3) grid (#2,#4);
  \draw[->, thick, charcoal] (#1,0)--(#2,0) node[below right, font=\tiny]{$x$};
  \draw[->, thick, charcoal] (0,#3)--(0,#4) node[left, font=\tiny]{$y$};}
\newcommand{\dpt}[2]{\fill[navy] (#1,#2) circle (2.6pt);}

\begin{document}

% TITLE SLIDE (CourseName is not defined in beamer, so the name is literal)
{
\setbeamercolor{background canvas}{bg=forest}
\begin{frame}[plain]
  \vspace{2.8cm}\hspace{0.55cm}%
  \begin{minipage}{0.9\textwidth}
    {\color{goldacc}\bfseries\small LESSON 2.5}\\[0.5cm]
    {\color{white}\bfseries\LARGE Linear Regression}\\[1.2cm]
    {\color{forestmid}\small Algebra 2: Shepherd~$\cdot$~Unit 2: Linear Functions}
  \end{minipage}
\end{frame}
}

% HOOK
\begin{frame}[t, plain]
  \forestheader{When the data won't line up}
  \sectionlabel{Hook}
  \begin{columns}[T]
    \begin{column}{0.54\textwidth}
      For each player: \textbf{hours practiced} ($x$) vs.\ \textbf{free throws made} out of $15$ ($y$).
      \begin{itemize}\setlength{\itemsep}{3pt}
        \item The points do \emph{not} lie on one line\ldots
        \item \ldots but they clearly \emph{trend} upward.
      \end{itemize}
      \vspace{4pt}
      No exact rule fits messy data --- but a \textbf{line of best fit} can summarize the trend.
    \end{column}
    \begin{column}{0.44\textwidth}
      \centering
      \begin{tikzpicture}[scale=0.34]
        \lgrid{-0.5}{9}{-0.5}{15}
        \dpt{1}{4}\dpt{2}{5}\dpt{3}{6}\dpt{4}{8}\dpt{5}{9}\dpt{6}{11}\dpt{7}{12}\dpt{8}{13}
      \end{tikzpicture}
    \end{column}
  \end{columns}
\end{frame}

% ASSOCIATION
\begin{frame}[t, plain]
  \forestheader{Describe it: direction, form, strength}
  \sectionlabel{Scatterplots}
  \begin{columns}[T]
    \begin{column}{0.34\textwidth}
      Three separate questions:
      \begin{itemize}\setlength{\itemsep}{4pt}
        \item \textbf{Direction:} up (positive), down (negative), or none?
        \item \textbf{Form:} does a \emph{line} fit?
        \item \textbf{Strength:} how tightly?
      \end{itemize}
    \end{column}
    \begin{column}{0.64\textwidth}
      \centering
      \begin{tikzpicture}[scale=0.3]
        \lgrid{-0.5}{7}{-0.5}{7}
        \dpt{1}{2}\dpt{2}{2}\dpt{3}{4}\dpt{4}{4}\dpt{5}{6}\dpt{6}{6}
        \node[font=\scriptsize] at (3,-1.7){positive};
      \end{tikzpicture}
      \hspace{0.2cm}
      \begin{tikzpicture}[scale=0.3]
        \lgrid{-0.5}{7}{-0.5}{7}
        \dpt{1}{6}\dpt{2}{5}\dpt{3}{5}\dpt{4}{3}\dpt{5}{2}\dpt{6}{2}
        \node[font=\scriptsize] at (3,-1.7){negative};
      \end{tikzpicture}
      \hspace{0.2cm}
      \begin{tikzpicture}[scale=0.3]
        \lgrid{-0.5}{7}{-0.5}{7}
        \dpt{1}{3}\dpt{2}{6}\dpt{3}{2}\dpt{4}{5}\dpt{5}{3}\dpt{6}{5}
        \node[font=\scriptsize] at (3,-1.7){none};
      \end{tikzpicture}
    \end{column}
  \end{columns}
\end{frame}

% CORRELATION COEFFICIENT
\begin{frame}[t, plain]
  \forestheader{One number: the correlation coefficient $r$}
  \sectionlabel[goldacc]{Goodness of fit}
  \begin{itemize}\setlength{\itemsep}{4pt}
    \item $r$ always lies in $[-1,1]$: \textbf{sign} $=$ direction, \textbf{size} $|r|$ $=$ strength.
    \item $r\approx +1$: tight rising line. \quad $r\approx -1$: tight falling line. \quad $r\approx 0$: no linear trend.
  \end{itemize}
  \vspace{2pt}
  \begin{center}
  \begin{tikzpicture}[scale=1]
    \draw[semithick, charcoal, ->] (-4.2,0)--(4.2,0);
    \foreach \x/\lab in {-4/-1, 0/0, 4/1}{\draw (\x,0.12)--(\x,-0.12) node[below, font=\scriptsize]{$\lab$};}
    \node[font=\scriptsize, forest] at (-4,0.45){strong negative};
    \node[font=\scriptsize, forest] at (0,0.45){no linear trend};
    \node[font=\scriptsize, forest] at (4,0.45){strong positive};
  \end{tikzpicture}
  \end{center}
  \begin{block}{Watch out}
    A \emph{negative} $r$ is not ``weak.'' $r=-0.9$ is a \textbf{strong} association --- and $r$ measures a
    \emph{linear} fit only.
  \end{block}
\end{frame}

% LINE OF BEST FIT
\begin{frame}[t, plain]
  \forestheader{The line of best fit: $\hat y = ax + b$}
  \sectionlabel{Regression}
  \begin{columns}[T]
    \begin{column}{0.54\textwidth}
      Technology returns \; $\hat y = 1.5x + 2$, \; $r=0.96$.
      \begin{itemize}\setlength{\itemsep}{4pt}
        \item \textbf{Slope} $1.5$: each extra hour of practice $\Rightarrow$ about $1.5$ more made (a per-unit rate).
        \item \textbf{Intercept} $2$: baseline made at $0$ hours.
        \item $\hat y$ is a \emph{predicted} value.
      \end{itemize}
    \end{column}
    \begin{column}{0.44\textwidth}
      \centering
      \begin{tikzpicture}[scale=0.34]
        \lgrid{-0.5}{9}{-0.5}{15}
        \draw[very thick, forest] (0,2)--(8.5,14.75);
        \dpt{1}{4}\dpt{2}{5}\dpt{3}{6}\dpt{4}{8}\dpt{5}{9}\dpt{6}{11}\dpt{7}{12}\dpt{8}{13}
      \end{tikzpicture}
    \end{column}
  \end{columns}
\end{frame}

% PREDICTING
\begin{frame}[t, plain]
  \forestheader{Predicting --- inside vs.\ outside the data}
  \sectionlabel[goldacc]{Interpolation / extrapolation}
  Substitute an $x$ into $\hat y = 1.5x+2$:
  \begin{itemize}\setlength{\itemsep}{5pt}
    \item $x=6$ (inside the data, $1$--$8$): $\hat y = 11$ --- \textbf{interpolation} (safer).
    \item $x=12$ (outside): $\hat y = 20$\ldots\ but only $15$ shots exist! \textbf{Extrapolation} can break down.
  \end{itemize}
  \vspace{2pt}
  \begin{block}{Always ask}
    Is this prediction \emph{reasonable}? A line predicting $20$ made out of $15$ is a red flag that you have
    extrapolated too far.
  \end{block}
\end{frame}

% CAUSATION + ROADMAP
\begin{frame}[t, plain]
  \forestheader{Correlation is not causation --- and where we go next}
  \sectionlabel{Connections}
  \textbf{Correlation $\ne$ causation.} Ice-cream sales and pool drownings both rise in summer ($r>0$), but ice
  cream does not \emph{cause} drownings --- \emph{hot weather} is the lurking variable. A strong $r$ means the
  variables move \emph{together}, not that one \emph{causes} the other.\\[8pt]
  \begin{itemize}\setlength{\itemsep}{4pt}
    \item \textbf{Today:} scatterplots \& association; the coefficient $r$; the line of best fit; predicting; causation.
    \item \textbf{Next:} the \textbf{Unit 2 Test}, then \textbf{Unit 3 --- Quadratic Functions}, where the curve of best fit is a \emph{parabola}.
  \end{itemize}
  \vspace{2pt}
  {\footnotesize Finding $r$ and the equation is a calculator/Desmos task; reading and interpreting them is on you.}
\end{frame}

\end{document}
